% !TEX program = xelatex
\documentclass[12pt,a4paper]{article}

% ---------- Font và ngôn ngữ ----------
\usepackage{fontspec}
\setmainfont{texgyretermes}[
  Extension      = .otf,
  UprightFont    = *-regular,
  BoldFont       = *-bold,
  ItalicFont     = *-italic,
  BoldItalicFont = *-bolditalic]
\setmonofont{DejaVu Sans Mono}[Scale=0.85]
\usepackage{polyglossia}
\setmainlanguage{vietnamese}

% ---------- Bố cục ----------
\usepackage[a4paper,
  left=3cm, right=3cm, top=2.5cm, bottom=2.5cm,
  headheight=15pt, headsep=0.8cm, footskip=1cm]{geometry}
\usepackage{graphicx}
\usepackage{etoolbox}
\usepackage{needspace}
\usepackage{amsmath}
\usepackage{setspace}
\setstretch{1.15}
\usepackage{parskip}
\usepackage{enumitem}
\setlist{nosep, leftmargin=*}

% ---------- Bảng, màu, hình ----------
\usepackage{array, booktabs, tabularx, longtable, multirow, makecell}
\usepackage[table]{xcolor}
\usepackage{tikz}
\usetikzlibrary{positioning, arrows.meta, fit, backgrounds, shapes.geometric, shapes.misc}
\usepackage{caption}
\captionsetup{font=small, labelfont=bf}
\usepackage{tcolorbox}
\tcbuselibrary{skins}

% ---------- Tiêu đề mục, header ----------
\definecolor{primary}{HTML}{7A1E3A}
\definecolor{accent}{HTML}{1F5FA8}
\definecolor{softgray}{HTML}{F3F3F3}
\usepackage{titlesec}
\titleformat{\section}{\Large\bfseries\color{primary}}{\thesection.}{0.5em}{}
\titleformat{\subsection}{\large\bfseries\color{accent}}{\thesubsection}{0.5em}{}
\titlespacing*{\section}{0pt}{1.2em}{0.5em}
\titlespacing*{\subsection}{0pt}{0.9em}{0.3em}

%% ------------------------------------------------------------
%% THÔNG TIN BÁO CÁO (dùng cho bìa, header, footer)
%% ------------------------------------------------------------
\newcommand{\tenDHQG}{ĐẠI HỌC QUỐC GIA THÀNH PHỐ HỒ CHÍ MINH}
\newcommand{\tenTruong}{TRƯỜNG ĐẠI HỌC KHOA HỌC TỰ NHIÊN}
\newcommand{\tenKhoa}{KHOA CÔNG NGHỆ THÔNG TIN}
\newcommand{\loaiBaoCao}{ĐỀ XUẤT ĐỒ ÁN CUỐI KỲ}
\newcommand{\tenNganh}{Môn: Vật lý cho Công nghệ thông tin}
\newcommand{\tenDeTaiChinh}{Phòng học thông minh}
\newcommand{\tenDeTaiPhu}{Giám sát chất lượng không khí và cảnh báo nhiệt độ, độ ẩm}
\newcommand{\tenDeTaiAnh}{Smart Classroom: Air Quality Monitoring and Temperature--Humidity Alert}
\newcommand{\tenDeTaiNgan}{Đề xuất đồ án Vật lý cho CNTT}   % hiển thị trên header
\newcommand{\giangVien}{Bùi Thanh Lâm, Cao Xuân Nam, Đặng Hoài Thương}
\newcommand{\tenNhom}{}   % để trống thì bìa không hiện dòng tên nhóm

\newcommand{\svA}{Nguyễn Gia Quốc Uy}  \newcommand{\mssvA}{24127261}
\newcommand{\svB}{Ngô Quốc Tuấn}       \newcommand{\mssvB}{24127581}
\newcommand{\svC}{Nguyễn Quốc Thành}   \newcommand{\mssvC}{24127542}

\newcommand{\khoaHoc}{2026 -- 2027}
\newcommand{\thangNam}{Tháng 10 năm 2026}
\newcommand{\diaDiem}{TP. Hồ Chí Minh}
\newcommand{\fileLogo}{logo-hcmus.png}   % đặt logo cùng thư mục (tùy chọn)

%% ------------------------------------------------------------
%% HEADER & FOOTER (theo format.tex)
%% ------------------------------------------------------------
\usepackage{fancyhdr}
\pagestyle{fancy}
\fancyhf{}
\renewcommand{\headrulewidth}{0.4pt}
\renewcommand{\footrulewidth}{0.4pt}

\fancyhead[L]{\small\textit{Trường ĐH Khoa học Tự nhiên -- ĐHQG-HCM}}
\fancyhead[R]{\small\textit{\tenDeTaiNgan}}
\fancyfoot[L]{\small\textit{\MakeUppercase{\tenKhoa}}}
\fancyfoot[R]{\small\thepage}

\fancypagestyle{plain}{%
  \fancyhf{}
  \renewcommand{\headrulewidth}{0pt}
  \renewcommand{\footrulewidth}{0.4pt}
  \fancyfoot[L]{\small\textit{\MakeUppercase{\tenKhoa}}}
  \fancyfoot[R]{\small\thepage}
}

%% ------------------------------------------------------------
%% TRANG BÌA (theo format.tex: khung viền kép)
%% ------------------------------------------------------------
\newcommand{\makecover}{%
  \begin{titlepage}
  \thispagestyle{empty}
  \singlespacing
  \begin{tikzpicture}[remember picture, overlay, inner sep=0, outer sep=0]
    % Khung chinh
    \draw[blue!55!black, line width=5.5pt]
      ([xshift=-1.5cm, yshift=-2cm]current page.north east) coordinate (A) --
      ([xshift=1.5cm, yshift=-2cm]current page.north west) coordinate (B) --
      ([xshift=1.5cm, yshift=2cm]current page.south west) coordinate (C) --
      ([xshift=-1.5cm, yshift=2cm]current page.south east) coordinate (D) -- cycle;

    % Vien long 1
    \draw
      ([yshift=0.2cm,xshift=-0.3cm]A)-- ([yshift=0.2cm,xshift=0.3cm]B)--
      ([yshift=-0.3cm,xshift=0.3cm]B) --([yshift=-0.3cm,xshift=-0.2cm]B)--
      ([yshift=0.3cm,xshift=-0.2cm]C)--([yshift=0.3cm,xshift=0.3cm]C)--
      ([yshift=-0.2cm,xshift=0.3cm]C)-- ([yshift=-0.2cm,xshift=-0.3cm]D)--
      ([yshift=0.3cm,xshift=-0.3cm]D)--([yshift=0.3cm,xshift=0.2cm]D)--
      ([yshift=-0.3cm,xshift=0.2cm]A)--([yshift=-0.3cm,xshift=-0.3cm]A)--
      ([yshift=0.2cm,xshift=-0.3cm]A);

    % Vien long 2
    \draw
      ([yshift=-0.2cm,xshift=0.2cm]A)-- ([yshift=-0.2cm,xshift=-0.2cm]B)--
      ([yshift=0.2cm,xshift=-0.2cm]B) --([yshift=0.2cm,xshift=0.2cm]B)--
      ([yshift=-0.2cm,xshift=0.2cm]C)--([yshift=-0.2cm,xshift=-0.2cm]C)--
      ([yshift=0.2cm,xshift=-0.2cm]C)-- ([yshift=0.2cm,xshift=0.2cm]D)--
      ([yshift=-0.2cm,xshift=0.2cm]D)--([yshift=-0.2cm,xshift=-0.2cm]D)--
      ([yshift=0.2cm,xshift=-0.2cm]A)--([yshift=0.2cm,xshift=0.2cm]A)--
      ([yshift=-0.2cm,xshift=0.2cm]A);
  \end{tikzpicture}
  \centering
  \vspace*{0.3cm}

  {\fontsize{14pt}{17pt}\selectfont\bfseries \tenDHQG\par}
  \vspace{2pt}
  {\fontsize{14pt}{17pt}\selectfont\bfseries \tenTruong\par}
  \vspace{2pt}
  {\fontsize{14pt}{17pt}\selectfont\bfseries \tenKhoa\par}
  \vspace{0.6cm}
  \rule{4cm}{0.8pt}

  \vspace{0.8cm}
  \IfFileExists{\fileLogo}{\includegraphics[height=3.5cm]{\fileLogo}}%
    {\fbox{\parbox[c][3.5cm][c]{3.5cm}{\centering\small [Logo HCMUS]\\(đặt file \texttt{\fileLogo})}}}

  \vspace{0.9cm}
  {\fontsize{14pt}{18pt}\selectfont\bfseries\color{black} \loaiBaoCao\par}
  \ifdefempty{\tenNganh}{}{\vspace{6pt}{\fontsize{15pt}{20pt}\selectfont\color{black} \tenNganh\par}}

  \vspace{0.75cm}
  \begin{minipage}{0.94\textwidth}
    \centering
    {\fontsize{19pt}{25pt}\selectfont\bfseries\color{black}
      \tenDeTaiChinh\par}
    \vspace{3pt}
    {\fontsize{17pt}{23pt}\selectfont\color{black}
      \tenDeTaiPhu\par}
    \vspace{0.3cm}
    {\fontsize{15pt}{20pt}\selectfont\itshape\color{black}
      (\tenDeTaiAnh)\par}
  \end{minipage}\par
  \ifdefempty{\tenNhom}{}{%
    \vspace{0.3cm}
    {\color{black}\rule{2.4cm}{0.5pt}\par}
    \vspace{0.25cm}
    {\fontsize{14pt}{18pt}\selectfont\bfseries\color{black} \tenNhom\par}}

  \vfill

  \begin{center}
  \begin{tabular}{@{}r@{\hspace{0.4cm}}l@{}}
    \textbf{Giảng viên hướng dẫn:} & Bùi Thanh Lâm \\[2pt]
                                   & Cao Xuân Nam \\[2pt]
                                   & Đặng Hoài Thương \\[4pt]
    \textbf{Sinh viên thực hiện:}  & \svA \ -- \mssvA \\[2pt]
                                   & \svB \ -- \mssvB \\[2pt]
                                   & \svC \ -- \mssvC \\[4pt]
    \textbf{Năm học:}              & \khoaHoc \\
  \end{tabular}
  \end{center}

  \vfill

  {\fontsize{13pt}{16pt}\selectfont\bfseries \diaDiem, \thangNam\par}
  \vspace{0.3cm}
  \end{titlepage}
  \setcounter{page}{1}
}

\usepackage[hidelinks]{hyperref}

\newcolumntype{L}[1]{>{\raggedright\arraybackslash}p{#1}}
\newcolumntype{C}[1]{>{\centering\arraybackslash}p{#1}}
\newcolumntype{Y}{>{\raggedright\arraybackslash}X}
% Bảng dữ liệu: viền ngoài đậm (I), đường kẻ bên trong nhạt (i)
\definecolor{gridgray}{gray}{0.7}
\arrayrulecolor{gridgray}
\newcolumntype{I}{!{{\color{black}\vrule width 0.9pt}}}
\newcolumntype{i}{!{{\color{gridgray}\vrule width 0.4pt}}}
\newcommand{\outerline}{\Xhline{0.9pt}}
\renewcommand{\arraystretch}{1.25}

\newtcolorbox{notebox}[1][]{colback=softgray, colframe=primary, boxrule=0.6pt,
  arc=2pt, left=6pt, right=6pt, top=4pt, bottom=4pt, #1}

\begin{document}

% ======================================================================
% TRANG BÌA
% ======================================================================
\makecover

\pagenumbering{roman}
\tableofcontents
\clearpage
\pagenumbering{arabic}

\noindent\begin{minipage}{\linewidth}
\setlength{\parskip}{6pt}
% ======================================================================
\section{Giới thiệu đề tài}
% ======================================================================
Phòng học thường đóng kín và đông người. Khi thông gió kém, nồng độ CO\textsubscript{2} tăng nhanh gây buồn ngủ, giảm tập trung, nhưng người trong phòng không tự nhận biết được. Nhiệt độ quá cao ảnh hưởng sức khỏe và làm thiết bị điện quá nhiệt; độ ẩm quá cao gây đọng hơi nước, độ ẩm quá thấp gây tĩnh điện.

Nhóm đề xuất một thiết bị IoT đặt trong phòng học, gồm board ESP32 đọc cảm biến nhiệt độ, độ ẩm (DHT22) và CO\textsubscript{2} ước lượng (MQ-135), gửi dữ liệu qua MQTT đến website do nhóm tự xây dựng. Website hiển thị số liệu theo thời gian thực và biểu đồ lịch sử, điều khiển đèn LED RGB báo mức chất lượng không khí, điều khiển còi cảnh báo khi nhiệt độ, độ ẩm ở mức nguy hiểm, và gửi thông báo Telegram. Sản phẩm được làm thật, đóng hộp, và người dùng không có kiến thức công nghệ thông tin vẫn tự cài đặt được qua giao diện.

\vspace{0.5cm}
\begin{notebox}
\begin{itemize}
  \item \textbf{Nhóm:} Thành (DHT22), Uy (MQ-135), Tuấn (LED RGB, còi); mỗi người phụ trách một luồng căn bản với thiết bị riêng.
  \item \textbf{Hình thức:} sản phẩm thật, đóng hộp.
  \item \textbf{Yêu cầu nâng cao dự kiến thực hiện:} \#3, \#4, \#5, \#6, \#9, \#10, \#12.
  \item \textbf{Chi phí linh kiện dự kiến:} khoảng 480.000đ.
\end{itemize}
\end{notebox}
\end{minipage}
\par\bigskip

\noindent\begin{minipage}{\linewidth}
\setlength{\parskip}{6pt}
% ======================================================================
\section{Phạm vi chức năng}
% ======================================================================
\begin{tabularx}{\linewidth}{I C{1.2cm} i Y i L{3.3cm} i C{1.7cm} I}
\outerline
\textbf{Mã} & \textbf{Use case} & \textbf{Yêu cầu tương ứng} & \textbf{Phụ trách} \\
\hline
UC1 & Xem nhiệt độ, độ ẩm theo thời gian thực & Căn bản, luồng (1) & Thành \\
\hline
UC2 & Xem CO\textsubscript{2} ước lượng theo thời gian thực & Căn bản, luồng (1) & Uy \\
\hline
UC3 & Điều khiển LED RGB từ web: tự động theo mức chất lượng không khí hoặc chọn màu thủ công & Căn bản, luồng (2) & Tuấn \\
\hline
UC4 & Còi cảnh báo khi nhiệt độ, độ ẩm vượt ngưỡng; thử và tắt còi từ web & Nâng cao \#3 & Tuấn \\
\hline
UC5 & Lưu dữ liệu cảm biến theo thời gian lên cloud & Nâng cao \#4 & Uy \\
\hline
UC6 & Xem biểu đồ lịch sử & Nâng cao \#5 & Thành \\
\hline
UC7 & Nhận thông báo Telegram khi chất lượng không khí xấu & Nâng cao \#6 & Tuấn \\
\hline
UC8 & Đăng ký, đăng nhập; tài khoản lưu trong database & Nâng cao \#9 & Thành \\
\hline
UC9 & Cấu hình WiFi qua giao diện khi cài đặt lần đầu & Nâng cao \#12 & Uy \\
\hline
UC10 & Xem mức chất lượng không khí (Tốt / Trung bình / Kém) & Phục vụ UC3, UC7 & Tuấn \\
\hline
UC11 & Chỉnh ngưỡng cảnh báo trên trang Cài đặt & Phục vụ UC4, UC7, UC10 & Uy \\
\outerline
\end{tabularx}

\vspace{0.5cm}
\noindent\textbf{Website do nhóm tự xây dựng, không dùng NodeRED (yêu cầu nâng cao \#10).}
\end{minipage}
\par\bigskip

\noindent\begin{minipage}{\linewidth}
\setlength{\parskip}{6pt}
% ======================================================================
\section{Kiến trúc hệ thống}
% ======================================================================
\subsection{Sơ đồ tổng thể}
\begin{center}
\centering
\resizebox{\linewidth}{!}{%
\begin{tikzpicture}[
  font=\small,
  >={Stealth[length=2.2mm]},
  dev/.style={draw, rounded corners=2pt, fill=#1, minimum width=2.5cm, minimum height=0.95cm, align=center},
  dev/.default=white,
  mod/.style={draw, fill=white, minimum width=2.8cm, minimum height=0.85cm, align=center, font=\footnotesize},
  svc/.style={draw, rounded corners=3pt, fill=accent!10, minimum width=3.2cm, minimum height=1.1cm, align=center},
  lbl/.style={font=\scriptsize, fill=white, inner sep=1.5pt},
]
% Thiết bị input / output
\node[dev=primary!12] (dht)   at (0, 3.0)  {DHT22\\\scriptsize nhiệt độ, độ ẩm (Thành)};
\node[dev=primary!12] (mq135) at (0, 1.6)  {MQ-135\\\scriptsize CO\textsubscript{2} ước lượng (Uy)};
\node[dev=accent!15]  (buz)   at (0, -0.6) {Buzzer\\\scriptsize còi cảnh báo (Tuấn)};
\node[dev=accent!15]  (led)   at (0, -2.2) {LED RGB WS2812\\\scriptsize (Tuấn)};

% Module firmware trong ESP32
\node[mod] (fdht)  at (4.4, 3.0)  {Đọc DHT22};
\node[mod] (fmq)   at (4.4, 1.6)  {Đọc MQ-135,\\quy đổi ppm};
\node[mod] (fbuz)  at (4.4, -0.6) {Nhận lệnh,\\bật/tắt còi};
\node[mod] (fled)  at (4.4, -2.2) {Nhận lệnh,\\điều khiển LED};
\begin{scope}[on background layer]
\node[draw=primary, thick, rounded corners=4pt, fill=primary!5, inner sep=9pt,
      fit=(fdht)(fmq)(fbuz)(fled), label={[font=\small\bfseries, text=primary]above:Board ESP32 (dùng chung)}] (esp) {};
\end{scope}

% Broker
\node[draw, ellipse, fill=softgray, minimum width=2.8cm, minimum height=1.4cm, align=center] (broker) at (8.8, 0.9) {MQTT Broker\\\scriptsize HiveMQ Cloud};

% Website
\node[svc] (be)   at (12.8, 2.4) {Backend\\\scriptsize Node.js, Express, mqtt.js};
\node[svc] (qual) at (15.6, 0.2) {Xử lý ngưỡng\\\scriptsize mức chất lượng, cảnh báo};
\node[svc] (fe)   at (17.6, 2.4) {Frontend\\\scriptsize đăng nhập, cài đặt\\\scriptsize text, gauge, chart\\\scriptsize nút LED, nút thử/tắt còi};
\begin{scope}[on background layer]
\node[draw=accent, thick, rounded corners=4pt, fill=accent!4, inner sep=9pt,
      fit=(be)(qual)(fe), label={[font=\small\bfseries, text=accent]above:Website tự xây dựng}] (web) {};
\end{scope}

\node[draw, cylinder, shape border rotate=90, aspect=0.25, fill=softgray, minimum width=2.6cm, minimum height=1.5cm, align=center] (fb) at (12.8, -3.0) {Firebase\\\scriptsize dữ liệu, tài khoản};
\node[draw, rounded corners=2pt, fill=white, minimum width=2.2cm, minimum height=0.95cm, align=center] (tg) at (15.6, -3.0) {Telegram\\\scriptsize Bot API};
\node[draw, rounded rectangle, fill=white, minimum width=2.2cm, minimum height=1cm, align=center] (user) at (18.4, -3.0) {Giảng viên /\\Sinh viên};

% Thiết bị <-> firmware
\draw[->] (dht) -- (fdht);
\draw[->] (mq135) -- (fmq);
\draw[->] (fbuz) -- (buz);
\draw[->] (fled) -- (led);

% Firmware <-> broker
\draw[->] (fdht.east) -- node[lbl, sloped, above] {topic dht22} (broker.north west);
\draw[->] (fmq.east) -- node[lbl, sloped, above] {topic mq135} (broker.west);
\draw[->] (broker.south) |- node[lbl, pos=0.75, above] {topic buzzer/cmd} (fbuz.east);
\draw[->] (broker.south) |- node[lbl, pos=0.75, above] {topic led/cmd} (fled.east);

% Broker <-> backend
\draw[->] (broker.north) |- node[lbl, pos=0.75, above] {dữ liệu cảm biến} ([yshift=6pt]be.west);
\draw[->] ([yshift=-8pt]be.west) -| node[lbl, pos=0.3, below] {lệnh LED, còi} (broker.north east);

% Trong website
\draw[<->] (be) -- node[lbl, above] {WebSocket / HTTP} (fe);
\draw[<->] (be.south east) -- node[lbl, sloped, above] {chế độ tự động} (qual.north west);

% Cloud, thông báo, người dùng
\draw[<->] (be.south) -- node[lbl, left, align=right] {ghi /\\đọc} (fb.north);
\draw[->] (qual.south) -- node[lbl, right] {cảnh báo} (tg.north);
\draw[<->] (fe.south -| user.north) -- (user.north);

% Cấu hình WiFi lần đầu
\node[draw, rounded rectangle, fill=white, minimum height=0.9cm, align=center] (phone) at (8.8, -4.4) {Điện thoại người dùng\\\scriptsize (lần đầu cài đặt)};
\draw[<->, dashed] (phone.west) -| node[lbl, pos=0.3, below] {WiFi qua captive portal} (esp.south);

% Chú thích
\node[dev=primary!12, minimum width=1.1cm, minimum height=0.5cm, font=\scriptsize] (k1) at (0, -4.4) {Input};
\node[dev=accent!15, minimum width=1.1cm, minimum height=0.5cm, font=\scriptsize, right=0.4cm of k1] (k2) {Output};
\end{tikzpicture}}
\captionof{figure}{Sơ đồ hoạt động của hệ thống.}
\label{fig:arch}
\end{center}
\end{minipage}
\par\bigskip

\subsection{Các luồng xử lý chi tiết}

\noindent\begin{minipage}{\linewidth}
\setlength{\parskip}{4pt}
\textbf{Luồng 1 --- Nhiệt độ, độ ẩm lên website} (UC1, Thành, căn bản luồng 1)
\begin{enumerate}[leftmargin=2em]
  \item DHT22 đo nhiệt độ và độ ẩm trong phòng.
  \item ESP32 đọc DHT22 vài giây một lần, đóng gói thành bản tin, ví dụ \texttt{\{"temp":28.5,"hum":62\}}.
  \item ESP32 gửi (publish) bản tin qua WiFi lên topic \texttt{dht22} tại HiveMQ Cloud.
  \item HiveMQ Cloud chuyển bản tin đến backend, vì backend đã đăng ký topic \texttt{dht22}.
  \item Backend kiểm tra dữ liệu hợp lệ và gắn thời điểm nhận.
  \item Backend đẩy số liệu đến trình duyệt qua WebSocket; frontend cập nhật ô text và gauge ngay, người dùng không cần tải lại trang.
\end{enumerate}
\end{minipage}
\par\medskip

\noindent\begin{minipage}{\linewidth}
\setlength{\parskip}{4pt}
\textbf{Luồng 2 --- CO\textsubscript{2} lên website} (UC2, Uy, căn bản luồng 1)
\begin{enumerate}[leftmargin=2em]
  \item MQ-135 thay đổi điện trở theo nồng độ khí, tạo ra một điện áp tương tự.
  \item ESP32 đọc điện áp qua ADC, tính điện trở cảm biến và quy đổi ra nồng độ CO\textsubscript{2} ước lượng (ppm) theo phương pháp ở mục~\ref{sec:mq135}.
  \item ESP32 publish bản tin lên topic \texttt{mq135} tại HiveMQ Cloud.
  \item HiveMQ Cloud chuyển bản tin đến backend.
  \item Backend kiểm tra dữ liệu hợp lệ và gắn thời điểm nhận.
  \item Backend đẩy số liệu đến frontend qua WebSocket; frontend cập nhật ô text và gauge.
\end{enumerate}
\end{minipage}
\par\medskip

\noindent\begin{minipage}{\linewidth}
\setlength{\parskip}{4pt}
\textbf{Luồng 3 --- Lưu trữ và biểu đồ lịch sử} (UC5 của Uy, UC6 của Thành)
\begin{enumerate}[leftmargin=2em]
  \item Mỗi khi nhận bản tin ở luồng 1 và luồng 2, backend đồng thời ghi giá trị kèm thời điểm nhận vào Firebase.
  \item Người dùng mở trang biểu đồ và chọn khoảng thời gian (theo giờ hoặc theo ngày).
  \item Frontend gửi yêu cầu HTTP đến backend.
  \item Backend truy vấn Firebase theo khoảng thời gian đã chọn và trả kết quả.
  \item Frontend vẽ biểu đồ nhiệt độ, độ ẩm, CO\textsubscript{2}.
\end{enumerate}
\end{minipage}
\par\medskip

\noindent\begin{minipage}{\linewidth}
\setlength{\parskip}{4pt}
\textbf{Luồng 4 --- Điều khiển LED RGB} (UC3, Tuấn, căn bản luồng 2)
\begin{enumerate}[leftmargin=2em]
  \item Người dùng chọn màu trên website, hoặc chuyển LED sang chế độ tự động.
  \item Frontend gửi lệnh đến backend.
  \item Backend kiểm tra người dùng đã đăng nhập. Ở chế độ tự động, backend quy đổi giá trị CO\textsubscript{2} mới nhất ra mức chất lượng (UC10) theo ngưỡng ở mục~\ref{sec:nguong} để chọn màu: xanh lá (Tốt), vàng (Trung bình), đỏ (Kém); mức chất lượng cũng được hiển thị trên website.
  \item Backend publish lệnh lên topic \texttt{led/cmd} tại HiveMQ Cloud.
  \item HiveMQ Cloud chuyển lệnh đến ESP32, vì ESP32 đã đăng ký topic \texttt{led/cmd}.
  \item ESP32 điều khiển LED WS2812 hiển thị màu tương ứng.
\end{enumerate}
Ở chế độ tự động, bước 3 đến bước 6 tự lặp lại mỗi khi luồng 2 làm mức chất lượng thay đổi.
\end{minipage}
\par\medskip

\noindent\begin{minipage}{\linewidth}
\setlength{\parskip}{4pt}
\textbf{Luồng 5 --- Còi cảnh báo} (UC4, Tuấn, nâng cao \#3)
\begin{enumerate}[leftmargin=2em]
  \item \emph{Thủ công:} người dùng bấm nút thử còi hoặc tắt còi; frontend gửi lệnh đến backend.
  \item \emph{Tự động:} mỗi khi luồng 1 mang số liệu mới, backend so nhiệt độ, độ ẩm với ngưỡng (mục~\ref{sec:nguong}); nếu vượt ngưỡng, backend tạo lệnh bật còi.
  \item Backend publish lệnh lên topic \texttt{buzzer/cmd}, ví dụ \texttt{\{"state":"on"\}}.
  \item HiveMQ Cloud chuyển lệnh đến ESP32.
  \item ESP32 bật hoặc tắt còi. Khi số liệu trở lại bình thường hoặc người dùng bấm tắt, backend gửi lệnh \texttt{\{"state":"off"\}}.
\end{enumerate}
\end{minipage}
\par\medskip

\noindent\begin{minipage}{\linewidth}
\setlength{\parskip}{4pt}
\textbf{Luồng 6 --- Thông báo Telegram} (UC7, Tuấn, nâng cao \#6)
\begin{enumerate}[leftmargin=2em]
  \item \emph{Liên kết một lần:} người dùng bấm nút ``Kết nối Telegram'', ứng dụng Telegram mở ra, bấm Start; bot báo cho backend mã cuộc trò chuyện và backend lưu mã này vào Firebase.
  \item \emph{Khi cảnh báo:} nếu CO\textsubscript{2} ở mức Kém liên tục quá thời gian quy định (mục~\ref{sec:nguong}), backend gọi Telegram Bot API.
  \item Telegram gửi tin nhắn cảnh báo đến điện thoại người dùng.
\end{enumerate}
\end{minipage}
\par\medskip

\noindent\begin{minipage}{\linewidth}
\setlength{\parskip}{4pt}
\textbf{Luồng 7 --- Đăng ký, đăng nhập} (UC8, Thành, nâng cao \#9)
\begin{enumerate}[leftmargin=2em]
  \item Người dùng nhập tài khoản trên frontend; frontend gửi đến backend.
  \item Backend đối chiếu với tài khoản lưu trong Firebase (khi đăng ký thì tạo mới).
  \item Nếu hợp lệ, backend cấp phiên đăng nhập; mọi thao tác xem và điều khiển trên website đều phải kèm phiên này.
  \item Lần đầu đăng nhập, người dùng nhập mã thiết bị để gắn thiết bị với tài khoản; backend lưu liên kết này vào Firebase.
\end{enumerate}
\end{minipage}
\par\medskip

\noindent\begin{minipage}{\linewidth}
\setlength{\parskip}{4pt}
\textbf{Luồng 8 --- Cấu hình WiFi lần đầu} (UC9, Uy, nâng cao \#12)
\begin{enumerate}[leftmargin=2em]
  \item Khi chưa có thông tin WiFi, ESP32 tự phát một mạng WiFi riêng.
  \item Người dùng kết nối bằng điện thoại; trang cấu hình tự hiện ra.
  \item Người dùng chọn WiFi nhà, nhập mật khẩu; ESP32 lưu vào bộ nhớ flash và khởi động lại.
  \item ESP32 kết nối WiFi, kết nối HiveMQ Cloud, rồi bắt đầu luồng 1 và luồng 2.
\end{enumerate}
\end{minipage}
\par\bigskip

\noindent\begin{minipage}{\linewidth}
\setlength{\parskip}{6pt}
\subsection{Ngưỡng mặc định}\label{sec:nguong}
Hệ thống có sẵn bộ ngưỡng dưới đây; người dùng có thể chỉnh lại trên trang Cài đặt (UC11).

\begin{center}
\begin{tabular}{I l i l i l I}
\outerline
\textbf{Đại lượng} & \textbf{Ngưỡng mặc định} & \textbf{Phản hồi của hệ thống} \\
\hline
\multirow{4}{*}{CO\textsubscript{2} ước lượng} & $< 800$ ppm & LED xanh lá (Tốt) \\
\cline{2-3}
 & $800$--$1200$ ppm & LED vàng (Trung bình) \\
\cline{2-3}
 & $> 1200$ ppm & LED đỏ (Kém) \\
\cline{2-3}
 & $> 1200$ ppm liên tục 10 phút & Gửi thông báo Telegram \\
\hline
Nhiệt độ & $> 35\,^{\circ}$C & Bật còi \\
\hline
Độ ẩm & $> 85\%$ hoặc $< 30\%$ & Bật còi \\
\outerline
\end{tabular}
\end{center}
\end{minipage}
\par\bigskip

\textbf{Ràng buộc thiết bị đôi một khác nhau:} DHT22 (Thành), MQ-135 (Uy), LED RGB (Tuấn) và Buzzer (Tuấn) mỗi thiết bị chỉ thuộc một luồng. Các yêu cầu nâng cao còn lại không dùng thêm thiết bị phần cứng.

\bigskip

\noindent\begin{minipage}{\linewidth}
\setlength{\parskip}{6pt}
% ======================================================================
\section{Cài đặt cho người dùng không chuyên}
% ======================================================================
Các cấu hình kỹ thuật (cloud, MQTT broker, Telegram bot) do nhóm thiết lập sẵn. Người dùng chỉ thao tác trên giao diện:
\begin{enumerate}
  \item \textbf{Kết nối WiFi (UC9):} lần đầu cắm điện, thiết bị phát một mạng WiFi riêng; người dùng kết nối bằng điện thoại, trang cấu hình tự hiện ra để chọn WiFi và nhập mật khẩu.
  \item \textbf{Tạo tài khoản (UC8):} đăng ký trên website và gắn thiết bị vào tài khoản.
  \item \textbf{Kết nối Telegram (UC7):} bấm một nút trên website, sau đó bấm \textbf{Start} trong Telegram.
  \item \textbf{Tùy chỉnh (UC11, không bắt buộc):} kéo thanh trượt để đổi ngưỡng cảnh báo; hệ thống có sẵn ngưỡng mặc định.
\end{enumerate}
\end{minipage}
\par\bigskip

\noindent\begin{minipage}{\linewidth}
\setlength{\parskip}{6pt}
% ======================================================================
\section{Công nghệ sử dụng}
% ======================================================================
\subsection{Thiết bị phần cứng}
\begin{tabularx}{\linewidth}{I L{2.3cm} i L{4.4cm} i Y I}
\outerline
\textbf{Thiết bị} & \textbf{Vai trò và cách kết nối} & \textbf{Nguyên lý vật lý} \\
\hline
ESP32 DevKit & Vi điều khiển trung tâm, lõi kép 240 MHz, có WiFi 2,4 GHz và ADC 12-bit; chạy firmware, kết nối WiFi và MQTT & Truyền dữ liệu bằng sóng điện từ 2,4 GHz; ADC lấy mẫu điện áp tương tự và chuyển thành số, đo tuyến tính tốt trong khoảng 0,15--2,45 V \\
\hline
DHT22 & Đo nhiệt độ ($-40$ đến $80\,^{\circ}$C) và độ ẩm (0--100\%); giao tiếp số qua 1 dây & Độ ẩm làm thay đổi hằng số điện môi của lớp polymer trong tụ điện, nên điện dung thay đổi; nhiệt độ làm thay đổi điện trở của nhiệt điện trở \\
\hline
MQ-135 & Đo chất lượng không khí, quy đổi ước lượng sang CO\textsubscript{2} tương đương; ngõ ra tương tự, qua cầu chia áp vào chân ADC & Lớp bán dẫn SnO\textsubscript{2} được nung nóng hấp phụ oxy trên bề mặt; khí khử (NH\textsubscript{3}, hơi cồn, benzen, khói) phản ứng với oxy này, trả electron về lớp bán dẫn và làm điện trở giảm. Cảm biến không chọn lọc riêng CO\textsubscript{2}, nên giá trị CO\textsubscript{2} chỉ là ước lượng \\
\hline
LED RGB WS2812 & Báo mức chất lượng không khí bằng màu; giao tiếp số 1 dây & Ba LED đỏ, lục, lam phát quang khi electron tái hợp với lỗ trống ở tiếp giáp p-n; màu cuối là tổng hợp cộng màu, độ sáng từng kênh điều chỉnh bằng PWM \\
\hline
Buzzer & Phát âm thanh cảnh báo; điều khiển qua chân GPIO và transistor & Hiệu ứng áp điện nghịch: điện áp biến thiên làm tấm gốm biến dạng, dao động và phát ra sóng âm \\
\outerline
\end{tabularx}
\end{minipage}
\par\bigskip

\noindent\begin{minipage}{\linewidth}
\setlength{\parskip}{6pt}
\subsection{Phương pháp ước lượng CO\textsubscript{2} từ MQ-135}\label{sec:mq135}
MQ-135 không đo trực tiếp nồng độ CO\textsubscript{2}; nhóm ước lượng qua bốn bước, nhằm cho giá trị ổn định và nhất quán chứ không đặt mục tiêu độ chính xác tuyệt đối.

\textbf{Bước 1 --- Từ điện áp ra điện trở cảm biến.} Cảm biến (điện trở $R_S$) mắc nối tiếp với điện trở tải $R_L$ trên module, cấp nguồn $V_C = 5$ V. Điện áp ngõ ra (tối đa 5 V) được đưa qua cầu chia áp, với $R_1$ nối tiếp và $R_2$ nối đất, sao cho điện áp vào ADC không vượt khoảng 2,5 V, nằm trong vùng đo tuyến tính của ADC ESP32:
\[
  V_{\text{out}} = V_{\text{ADC}}\,\frac{R_1 + R_2}{R_2},
  \qquad
  R_S = R_L\,\frac{V_C - V_{\text{out}}}{V_{\text{out}}}.
\]

\textbf{Bước 2 --- Hiệu chuẩn $R_0$.} Sau khi làm nóng cảm biến 24--48 giờ, đặt thiết bị ở ngoài trời thoáng, nơi nồng độ CO\textsubscript{2} xấp xỉ mức nền khí quyển $C_0 \approx 420$ ppm, đo $R_S$ và tính điện trở chuẩn
\[
  R_0 = R_S \left(\frac{C_0}{a}\right)^{1/b}.
\]
$R_0$ được lưu vào bộ nhớ flash của ESP32 và dùng cho các lần đo sau.

\textbf{Bước 3 --- Quy đổi theo đường đặc tuyến.} Đường đặc tuyến của nhà sản xuất cho CO\textsubscript{2} có dạng hàm lũy thừa:
\[
  C_{\text{CO}_2} = a\left(\frac{R_S}{R_0}\right)^{-b},
  \qquad a \approx 116{,}6,\; b \approx 2{,}769.
\]

\textbf{Bước 4 --- Lọc nhiễu.} Mỗi lần đo lấy trung bình nhiều mẫu ADC liên tiếp, đọc điện áp đã hiệu chỉnh của ESP32 (ADC của ESP32 không tuyến tính ở hai đầu thang đo), sau đó lấy trung bình trượt theo thời gian để loại các dao động ngắn.

\textbf{Giới hạn.} Vì SnO\textsubscript{2} phản ứng với nhiều loại khí khử, sự xuất hiện của hơi cồn, khói hoặc NH\textsubscript{3} sẽ làm giá trị ước lượng tăng theo. Vì vậy hệ thống dùng giá trị này để phân mức chất lượng không khí (Tốt / Trung bình / Kém) thay vì công bố như một phép đo CO\textsubscript{2} chính xác.
\end{minipage}
\par\bigskip

\noindent\begin{minipage}{\linewidth}
\setlength{\parskip}{6pt}
\subsection{Công nghệ phần mềm}
\begin{tabularx}{\linewidth}{I L{2.6cm} i L{5.0cm} i Y I}
\outerline
\textbf{Thành phần} & \textbf{Công nghệ} & \textbf{Vai trò} \\
\hline
Firmware & C++ trên Arduino framework, PlatformIO; thư viện DHT, Adafruit NeoPixel, PubSubClient, ArduinoJson & Đọc cảm biến, điều khiển LED và còi, gửi/nhận bản tin MQTT \\
\hline
Cấu hình WiFi & Thư viện WiFiManager & Phát WiFi riêng và trang cấu hình khi cài đặt lần đầu \\
\hline
Giao thức & MQTT (publish/subscribe) qua TLS & Truyền dữ liệu giữa ESP32 và backend theo yêu cầu đồ án \\
\hline
MQTT broker & HiveMQ Cloud (gói miễn phí) & Trung chuyển bản tin; thiết bị và backend kết nối ra ngoài nên không cần cấu hình router \\
\hline
Backend & Node.js, Express, mqtt.js, WebSocket & Nhận dữ liệu, xử lý ngưỡng, cung cấp API cho website, đẩy dữ liệu thời gian thực \\
\hline
Máy chủ & Laptop của nhóm & Chạy backend và phục vụ website khi demo, vấn đáp; khi triển khai thực tế cần máy chủ chạy liên tục \\
\hline
Frontend & HTML, CSS, JavaScript; Chart.js; thư viện gauge & Giao diện đăng nhập, số liệu, biểu đồ, điều khiển, cài đặt \\
\hline
Cơ sở dữ liệu & Firebase (Firestore) & Lưu dữ liệu theo thời gian, tài khoản, ngưỡng cảnh báo \\
\hline
Thông báo & Telegram Bot API & Gửi cảnh báo đến điện thoại người dùng \\
\hline
Quản lý mã nguồn & Git, GitHub & Tích hợp code của cả nhóm trong một repository \\
\outerline
\end{tabularx}
\end{minipage}
\par\bigskip

\noindent\begin{minipage}{\linewidth}
\setlength{\parskip}{6pt}
% ======================================================================
\section{Chi phí dự kiến}
% ======================================================================
\begin{center}
\begin{tabular}{I L{8cm} i r I}
\outerline
\textbf{Hạng mục} & \textbf{Giá dự kiến} \\
\hline
ESP32 DevKit & 150.000đ \\
\hline
DHT22 & 100.000đ \\
\hline
MQ-135 & 50.000đ \\
\hline
LED RGB WS2812 & 40.000đ \\
\hline
Buzzer & 20.000đ \\
\hline
Adapter nguồn 5 V -- 2 A và cáp & 40.000đ \\
\hline
Board đồng lỗ, dây, điện trở, transistor; hộp nhựa có lỗ thông khí & 80.000đ \\
\hline
\textbf{Tổng} & \textbf{$\approx$ 480.000đ} \\
\outerline
\end{tabular}
\end{center}
Phần mềm và hạ tầng (HiveMQ Cloud, Firebase, Telegram Bot API) dùng gói miễn phí.
\end{minipage}
\par\bigskip

\end{document}
